\scanpage{4}
Structure à l'infini des $M_{g,\nu}$

1. \uline{Courbes \unsure{standard|standart}}.

Soit $k$ un corps alg. clos. Une ``courbe \unsure{standard|standart}'' sur $k$ est un\deleted{e \unsure{courbe}} schéma $X$ sur $k$ satisfaisant les conditions suivantes :

a) $X$ quasi-projectif, \inserted{toute composante irréd. est} de dim $1$

b) Tout point de $X$ est soit lisse, soit un ``point quadratique'' (``ordinaire'') -- i.e. isom. (loc. ét) à la courbe $\operatorname{Spec}(k[X,Y]/XY)$ au point $0$.

Il est connu qu'on peut donner une unique ``complétion'' $\widehat X$ de $X$, telle que \inserted{$\widehat X$ soit un schéma propre, \unsure{que}} $X$ s'identifie à un ouvert dense de $\widehat X$, et que $\widehat X$ soit lisse en les pts de $\widehat X\setminus X=I$. Alors $\widehat X$ est une courbe projective, $I$ est \deleted{\unsure{contenu}} une partie finie de $\widehat X(k)$ contenue dans l'ens. ouvert des pts de lissité de $\widehat X$. L'ens. $A$ des pts singuliers de $X$ s'identifie à l'ens. analogue pour $\widehat X$
\[
(1)\qquad
\left\{\begin{aligned}
I&=\widehat X\setminus X\subset\widehat X^{\mathrm{lisse}}(k)
&&\text{``pts à l'infini'' de }X\\
A&=X\setminus X^{\mathrm{lisse}}=\widehat X\setminus\widehat X^{\mathrm{lisse}}\subset X(k)
&&\text{(pts singuliers de }X\text{)}
\end{aligned}\right.
\]
La donnée de $X$ équivaut à celle de
\mspage{3}
